%%%%%%%%%%%%%%%%%%%%%%%%%%%%%%%%%%%%%%%%%%%%%%%%%%%%%%%%%%%%
%2345678901234567890123456789012345678901234567890123456789012345678901234567890
%        1         2         3         4         5         6         7         8

\documentclass[letterpaper, 10 pt, conference]{ieeeconf}  % Comment this line out
                                                          % if you need a4paper
%\documentclass[a4paper, 10pt, conference]{ieeeconf}      % Use this line for a4
                                                          % paper

\IEEEoverridecommandlockouts                              % This command is only
                                                          % needed if you want to
                                                          % use the \thanks command
\overrideIEEEmargins
% See the \addtolength command later in the file to balance the column lengths
% on the last page of the document



% The following packages can be found on http:\\www.ctan.org
%\usepackage{graphics} % for pdf, bitmapped graphics files
%\usepackage{epsfig} % for postscript graphics files
%\usepackage{mathptmx} % assumes new font selection scheme installed
%\usepackage{times} % assumes new font selection scheme installed
%\usepackage{amsmath} % assumes amsmath package installed
%\usepackage{amssymb}  % assumes amsmath package installed
\usepackage{graphicx}
\usepackage{verbatim}
\usepackage{multirow}
\usepackage{rotating}
\usepackage{moreverb}                    % for boxedboxedverbatim
\usepackage{array}
\usepackage{fancyvrb}
\usepackage{multicol}
\usepackage{mdwlist}
\usepackage{tikz-er2}


\newtheorem{defn}{Definition}

\newcommand{\class}[1] {\textit{#1}}
\newcommand{\const}[1] {$\mathit{#1}$}
\newcommand{\objvar}[1] {$\mathsf{#1}$}
\newcommand{\stvar}[1] {\textsf{#1}}
\newcommand{\op}[1] {\textsl{#1}}
\newcommand{\nil} {\textit{nil}\ }

\newcommand\T{\rule{0pt}{2.6ex}}
\newcommand\B{\rule[-1.2ex]{0pt}{0pt}}

\definecolor{violetred}{cmyk}{0,0.85,0.31,0.18}
\definecolor{darkblue}{cmyk}{1,1,0,0.45}
\definecolor{lavenderblush4}{cmyk}{0,0.06,0.04,0.45}
\definecolor{packergreen}{cmyk}{0.46,0,0.21,0.76}
\definecolor{fuschia}{cmyk}{0,100,0,0}
\definecolor{graycmyk}{cmyk}{0,0,0,0.74}

\usetikzlibrary{calc,trees,positioning,arrows,chains,shapes.geometric,%
    decorations.pathreplacing,decorations.pathmorphing,shapes,%
    matrix,shapes.symbols,positioning,shadows}

% styles for flowcharts

\tikzstyle{every entity} = [top color=white, bottom color=blue!30,
                            draw=blue!50!black!100, drop shadow]
\tikzstyle{empty} = [top color=white, bottom color=white,
                            draw=white]

\tikzstyle{every weak entity} = [drop shadow={shadow xshift=.7ex,
                                 shadow yshift=-.7ex}]
\tikzstyle{every attribute} = [top color=white, bottom color=green!20,
                               draw=green, node distance=1cm, drop shadow]
\tikzstyle{ELLIPSE} = [draw, ellipse, top color=white, bottom color=green!20, draw=green, drop shadow]
\tikzstyle{MainAttribute} = [draw, rectangle,top color=white, bottom color=red!20,
                               draw=red, node distance=1cm, drop shadow]
\tikzstyle{DATABASE} = [draw, rectangle, rounded corners,top color=white, bottom color=graycmyk!50, draw=graycmyk, inner sep=10pt, drop shadow={shadow xshift=.7ex, shadow yshift=-.7ex}]


\tikzstyle{myarrow}=[->, >=stealth', thick, shorten <=2pt,shorten >=2pt]



\tikzstyle{output} = [draw, rectangle, rounded corners,top color=white, bottom color=red!30, draw=red, inner sep=10pt, drop shadow={shadow xshift=.7ex, shadow yshift=-.7ex}]

\tikzstyle{abstract}=[rectangle, draw=black, rounded corners, fill=blue!40, drop shadow,
        text centered, anchor=north, text=white, text width=3cm]


\newbox{\LegendOutput}
\savebox{\LegendOutput}{
    (\begin{tikzpicture}[]
    \node[output] (2) {\hspace{5 mm}};
    \end{tikzpicture}
    )}


\newenvironment{mylisting}
{\begin{list}{}{\setlength{\leftmargin}{1em}}\item\small}
{\end{list}}

\newenvironment{mytinylisting}
{\begin{list}{}{\setlength{\leftmargin}{1em}}\item\tiny\bfseries}
{\end{list}}


\title{\LARGE \bf
Metrics and Test Methods for Industrial Kit Building
}

%\author{ \parbox{3 in}{\centering Stephen Balakirsky\\
%         Intelligent Systems Division\\
%         National Institute of Standards and Technology\\
%        Gaithersburg, MD 20899, USA\\
%         {\tt\small stephen.balakirsky@nist.gov}}
%         \hspace*{ 0.5 in}
%         \parbox{3 in}{ \centering Zeid Kootbally\\
%          Department of Mechanical Engineering \\
%         University of Maryland\\
%         College Park, MD 20742, USA\\
%         {\tt\small zeid.kootbally@nist.gov}}\\ \\
%	\parbox{2.25 in}{\centering Craig Schlonoff\\
%         Intelligent Systems Division\\
%         National Institute of Standards and Technology\\
%        Gaithersburg, MD 20899, USA\\
%         {\tt\small craig.schlenoff@nist.gov}}
%        \hspace*{0.05in}
%         \parbox{2.25 in}{ \centering Thomas Kramer\\
%          Department of Mechanical Engineering \\
%         Catholic University of America\\
%         Washington, DC 20064, USA\\
%         {\tt\small thomas.kramer@nist.gov}}
%        \hspace*{0.05in}
%         \parbox{2.25 in}{ \centering Satyandra K. Gupta\\
%          Maryland Robotics Center\\
%         University of Maryland\\
%         College Park, MD 20742, USA\\
%         {\tt\small skgupta@umd.edu}}\\
%}

\author{Stephen Balakirsky, Thomas Kramer, and Zeid Kootbally% <-this % stops a space
\thanks{S. Balakirsky is with the Intelligent Systems Division, National Institute of Standards and Technology, Gaithersburg, MD, USA (e-mail:stephen.balakirsky@nist.gov)}% <-this % stops a space
\thanks{Z. Kootbally is with the Department of Mechanical Engineering, University of Maryland, College Park, MD, USA (email: zeid.kootbally@nist.gov)}%
\thanks{T. Kramer is with the Department of Mechanical Engineering, Catholic University of America, Washington, DC, USA (email: thomas.kramer@nist.gov)}%
} 